\chapter{实验结果、对比与讨论}

\section{结果汇报结构}
对每个场景，本文按“最优方案 -- 成本分解 -- 稳定性”三层结构输出结果：
\begin{enumerate}
  \item 最优检查次数 $n^\star$ 与最优检查点序列 $(k_1^\star,\ldots,k_{n^\star}^\star)$；
  \item 单位合格品最小成本 $\phi^\star$ 及其构成；
  \item 参数扰动下的策略变化与成本敏感性。
\end{enumerate}

\section{场景 I：理想判定结果分析}
当 $R_1=1,\ R_2=0$ 时，不存在误判停机项，成本主要由“废品损失、检查费用、恢复费用、主动换刀费用”构成。一般可观察到：
\begin{itemize}
  \item 检查次数过少会显著增加故障延迟发现导致的废品损失；
  \item 检查次数过多会抬升检查成本且边际收益下降；
  \item 最优方案通常出现在中等检查频度区间。
\end{itemize}

\section{场景 II：存在误判结果分析}
当 $R_1=0.98,\ R_2=0.4$ 时，误判机制引入“过早停机”风险。该场景下的最优策略往往与场景 I 有两点差异：
\begin{enumerate}
  \item 检查节点可能适当前移，以减少故障后长区间损失；
  \item 总检查次数不一定增加，需兼顾误停带来的额外代价。
\end{enumerate}
因此“更密集检查并不必然更优”，最优策略由误判成本与故障损失共同决定。

\section{成本分解示例表}
\begin{table}[H]
\centering
\caption{最优方案成本分解（示例模板，提交时填入实算值）}
\begin{tabular}{lcc}
\toprule
成本项 & 场景 I / 元 & 场景 II / 元 \\
\midrule
废品损失 $\theta_1\zeta_3$ & -- & -- \\
检查费用 $\theta_2\gamma$ & -- & -- \\
误停损失 $\theta_5\eta_1$ & 0 & -- \\
主动换刀 $\theta_4\eta_2$ & -- & -- \\
恢复费用 $\theta_3(1-\eta_1-\eta_2)$ & -- & -- \\
\midrule
总损失 $h$ & -- & -- \\
合格品期望 $\zeta_1$ & -- & -- \\
单位成本 $\phi=h/\zeta_1$ & -- & -- \\
\bottomrule
\end{tabular}
\end{table}

\section{灵敏度分析}
围绕关键参数 $\theta_1,\theta_2,\theta_3,\theta_5$ 做扰动实验，可得到如下规律：
\begin{itemize}
  \item $\theta_1$（废品损失）增大时，最优策略趋向于“更早检查”；
  \item $\theta_2$（检查成本）增大时，最优检查次数下降；
  \item $\theta_5$（误停损失）增大时，双检确认策略更具优势；
  \item $\theta_3$（恢复费用）增大时，预防性换刀倾向增强。
\end{itemize}

\section{改进策略讨论}
在第二问场景中，可引入“首次不合格即复检”的双检机制。其效果取决于以下条件：
\[
\text{误停代价较高且复检成本较低} \Longrightarrow \text{双检更优}.
\]
若检查本身代价较高或复检准确率提升有限，则双检机制未必优于单检策略。

\section{本章小结}
本章给出了两类场景下的结果解读框架与敏感性分析逻辑，指出最优策略的本质是“多成本项权衡”而非单一参数驱动。
